\documentclass[a4paper,11pt]{article} %tipo de documento y tamaño de papel y letra
\usepackage[utf8]{inputenc} %codificación de caracteres
\usepackage[spanish]{babel} %idioma
\usepackage{fancyhdr} %tipo de página, LINDA biggrin.gif
\usepackage[top=3.5cm,bottom=2.5cm,right=2cm,left=2.5cm]{geometry} %márgenes
\usepackage[rflt]{floatflt} %ni puta idea


%\usepackage{helvet} %esto es pa escribir con Arial en vez de times new roman
%\renewcommand\familydefault{\sfdefault} %descomenta estas lineas para escribir en arial

\usepackage{graphicx} %para usar imagenes
%\newcommand{\imgdir}{doc-img} % para meter las imágenes en una carpeta especial, tunz en la carpeta del documento tiene q ir otra que se llame 'doc-img'
\graphicspath{{\pic/}} %le dice que las imágenes están en la carpeta de arriba

\usepackage{amsmath} %pa usar símbolos matematicos
\numberwithin{equation}{section} % pa usar ecuaciones de modo lindo
\numberwithin{figure}{section} %para agregar imágenes
\numberwithin{table}{section} %para agregar tablas

\usepackage{subfigure} % pa usar sub figuras



\pagestyle{fancy} %configuración para la pagina linda
\renewcommand{\sectionmark}[1]{\markboth{}{\thesection\ \ #1}} %cambios de comentarios
\lhead{} %parte de arriba, izq
\chead{} %parte de arriba, centro
\rhead{\rightmark} %parte de arriba, derecha, le agrega la marca del capitulo
\lfoot{} %parte de abajo, izq
\cfoot{} %parte de abajo, centro
\rfoot{\thepage} %parte de abajo, derecha, va el numero de la página


%-------------portada---------------------------------%
\begin{document}
\begin{titlepage} %portada
\thispagestyle{empty} %borrar el formato de pagina linda
\begin{flushleft} %alinear a la izq
\includegraphics[scale=0.25]{logoUSMC.eps} %agregar imagen, en mi caso es el logo del dpto, pero tu podríai poner el de qca.
\vfill %vertical fill, llena hacia abajo
\end{flushleft}

\vspace{2cm} %espacio vertical , en realidad es un enter de 2 cm
\begin{center} %centrar
{\huge Laboratorio de análisis de señales: Experiencia 7 \\ %
Sistemas Lineales Invariantes en el Tiempo\\
Realizada el 3 de Diciembre de 2008
}
\end{center}

\vspace{6cm}
\begin{flushright} %alinear derecha
Nombre: Javier Palma E.\\
Curso: Laboratorio Análisis de Señales\\
Profesor: Alex Ulloa\\
10 de Diciembre 2008
\end{flushright}
\end{titlepage}
%------------------fin de la portada --------------------%

\setcounter{page}{1} %empezar enumerando la página 1

\tableofcontents %índice
\newpage

\section{} %Nombre_de_la_sección (capitulo)
\subsection{} %subsección, osea 1.1
\subsubsection{} %1.1.1
\subsubsubsection{} %1.1.1.1

\begin{figure}[h] %insertar imagen, la h significa 'here', osea aquí
\centering
\includegraphics[scale=0.6]{} % nombre de la imagen entre las llaves
\caption{} %leyenda bajo la imagen
\label{} %etiqueta para ubicar la imagen despues
\end{figure}

{\bf } %escribir en negrita
{\tt } %para escribir con letra tipo maquina de escribir, sirve para anotar códigos o en tu caso, fórmulas o reacciones químicas

\begin{itemize} %para enumerar, y usar viñetas tipo
%* hola
%* como estas
\item{} %este será el texto puntuado

\end{itemize}

\begin{table}[h] %tablas, h significa here
\begin{center}
\begin{tabular}{|c|c|} %la barrita representa la linea de separacion, y la 'c' la cantidad de columnas
\hline
primera columna & segunda columna \\ %el '&' es el separador de columnas
\hline
\\
\hline
\\
\hline
\\
\hline
\\ %y asi sucesivamente, no existe límite de filas
\end{tabular}
\end{center}
\caption{}
\label{}
\end{table}



\end{document}
